\section{余下账本事件导航：early-ming}

\label{sec:remaining-event-anchors-early-ming}

本节为尚未获得正文目标的明代账本记录提供跳转锚点。锚点标记可核日期与事件；朝廷、法律、家庭、城市、环境和海外网络的解释仍应由各条来源和已有叙事的材料界限约束。

\subsection{事件导航与来源边界}

\eventanchor{EM-1409-NURGAN}{1409（永乐七年）·明设奴儿干都司，尝试经黑龙江下游联络和安置东北诸部}

\eventanchor{EM-1410-FIRST-MONGOL-CAMPAIGN}{1410（永乐八年）·永乐帝亲征蒙古，北上击败本雅失里及相关部众}

\eventanchor{EM-1411-GRAND-CANAL}{1411（永乐九年）·朝廷大规模疏浚会通河并调整运河漕运工程}

\eventanchor{EM-1413-GUIZHOU-PROVINCE}{1413（永乐十一年）·明置贵州承宣布政使司，调整卫所与土司并存的西南治理}

\eventanchor{EM-1413-ZHENGHE-FOURTH-VOYAGE}{1413（永乐十一年）·郑和第四次下西洋，航程扩展至忽鲁谟斯等地}

\eventanchor{EM-1414-TULA}{1414（永乐十二年）·永乐帝北征瓦剌，在土拉河一带交战}

\eventanchor{EM-1415-CANAL-TRANSPORT}{1415（永乐十三年）·会通河航路投入漕粮北运，江南粮食可较稳定转输北京方向}

\eventanchor{EM-1416-BEIJING-CONSTRUCTION}{1416（永乐十四年）·北京宫城、坛庙和城防工程持续集中施工}

\eventanchor{EM-1417-ZHENGHE-FIFTH-VOYAGE}{1417（永乐十五年）·郑和第五次航行护送使节并继续西抵印度洋港口}

\eventanchor{EM-1420-BEIJING-CAPITAL}{1420（永乐十八年）·永乐帝正式定北京为京师并完成迁都制度安排}

\eventanchor{EM-1421-PALACE-FIRE}{1421（永乐十九年）五月·北京奉天、华盖、谨身三殿遭火灾}

\eventanchor{EM-1422-SECOND-NORTHERN-CAMPAIGN}{1422（永乐二十年）·永乐帝再次北征，寻击阿鲁台而未获决定性会战}

\eventanchor{EM-1423-THIRD-NORTHERN-CAMPAIGN}{1423（永乐二十一年）·永乐帝第三次北征并继续搜寻阿鲁台部众}

\eventanchor{EM-1424-HONGXI-REVERSALS}{1424（洪熙元年）·仁宗停止部分远征和营建项目，赦免、召还一批永乐朝获罪者}

\eventanchor{EM-1424-YONGLE-DEATH}{1424（永乐二十二年）七月·永乐帝北征归途中去世，太子朱高炽继位}

\eventanchor{EM-1425-XUANDE-ACCESSION}{1425（洪熙元年）五月·仁宗去世，皇太子朱瞻基即位为宣德帝}

\eventanchor{EM-1426-GAOXU-REBELLION}{1426（宣德元年）·汉王朱高煦在乐安起兵，宣德帝亲征后其降伏并被处置}

\eventanchor{EM-1427-JIAOZHI-WITHDRAWAL}{1427（宣德二年）·明军在交趾战局恶化后与黎利方面议和，准备撤军}

\eventanchor{EM-1428-LE-LOI-RECOGNITION}{1428（宣德三年）·明接受黎利政权的朝贡关系并调整安南外交称谓}

\eventanchor{EM-1430-YUQIAN-INSPECTION}{1430（宣德五年）·于谦奉命巡抚河南、山西，处理灾荒、军政与地方事务}

\eventanchor{EM-1430-ZHENGHE-SEVENTH-VOYAGE}{1430（宣德五年）·郑和第七次下西洋自南京出发}

\eventanchor{EM-1431-TIANFEI-STELE}{1431（宣德六年）·郑和等立《天妃灵应之记》碑，记述历次航行与祈祷}

\eventanchor{EM-1433-ZHENGHE-RETURN}{1433（宣德八年）·第七次航行舰队返抵中国，郑和卒年与卒地存在不同说法}

\eventanchor{EM-1435-XUANDE-DEATH}{1435（宣德十年）正月·宣德帝去世，九岁皇太子朱祁镇即位为英宗}

\eventanchor{EM-1436-LUCHUAN-CRISIS}{1436（正统元年）·麓川思任发扩张引发边境冲突，明廷开始部署征讨}

\eventanchor{EM-1438-FIRST-LUCHUAN-CAMPAIGN}{1438（正统三年）·王骥等统兵发动第一次大规模征麓川行动}

\eventanchor{EM-1440-ESEN-RAID}{1440（正统五年）·也先势力进犯明北边，边镇紧张加剧}

\eventanchor{EM-1441-SECOND-LUCHUAN-CAMPAIGN}{1441（正统六年）·明廷再遣大军征麓川，试图追捕思任发并重建土司秩序}

\eventanchor{EM-1442-ZHANG-TAIHOU-DEATH}{1442（正统七年）·张太皇太后去世，幼主时期的辅政平衡发生变化}

\eventanchor{EM-1443-THIRD-LUCHUAN-CAMPAIGN}{1443（正统八年）·明廷继续第三次征麓川并以招抚、分置应对残余势力}

\eventanchor{EM-1445-ESEN-TRIBUTE-CONFLICT}{1445（正统十年）·也先借朝贡、互市和使团规模问题向明廷施压，北边交涉恶化}

\eventanchor{EM-1448-DENG-MAOQI}{1448（正统十三年）·福建邓茂七起事扩大，地方官军与中央调兵镇压}

\eventanchor{EM-1448-THIRD-LUCHUAN-EXPEDITION}{1448（正统十三年）·明廷为麓川问题再度大规模调兵，征发延及多省军镇}

\eventanchor{EM-1449-BEIJING-DEFENSE}{1449（正统十四年）九月·于谦等组织北京守备，拒绝南迁并抵御瓦剌进逼}

\eventanchor{EM-1449-JINGTAI-ACCESSION}{1449（正统十四年）九月·郕王朱祁钰即帝位为景泰帝，尊英宗为太上皇}

\eventanchor{EM-1449-TUMU}{1449（正统十四年）八月·英宗亲征瓦剌，明军在土木堡溃败，英宗被俘}

\eventanchor{EM-1450-YINGZONG-RETURN}{1450（景泰元年）八月·瓦剌释放英宗，英宗返京后被安置为太上皇}

\eventanchor{ML-1409-001}{1409（永乐七年）一月·宝藏航次：第三次航次下达指令}

\eventanchor{ML-1409-002}{1409（永乐七年）二月15日·宝藏之旅：加勒三语铭文诞生}

\eventanchor{ML-1409-003}{1409（永乐七年）六月·卫拉特从明朝廷获得王子头衔}

\eventanchor{ML-1409-004}{1409（永乐七年）夏·宝藏之旅：宝船队返回中国}

\eventanchor{ML-1409-005}{1409（永乐七年）九月23日·克鲁伦之战：明军被乌列帖木儿汗击败}

\eventanchor{ML-1409-006}{1409（永乐七年）十月·郑和下西洋：率领27000人，按常规路线出发}

\eventanchor{ML-1409-007}{1409（永乐七年）十二月·中国第四次统治越南：明朝军队攻占陈鹗，但陈贵康成为叛军领袖}

\eventanchor{ML-1410-001}{1410（永乐八年）六月15日·第一次蒙古战役：永乐皇帝在斡难河畔击败额列帖木儿汗}

\eventanchor{ML-1410-002}{1410（永乐八年）七月·第一次蒙古战役：明军击败大兴安岭以东的阿鲁台并撤至南京}

\eventanchor{ML-1410-003}{1410（永乐八年）·明科特战争：宝藏船队登陆锡兰加勒并俘获甘波拉王国国王维杰亚巴胡六世}

\eventanchor{ML-1411-001}{1411（永乐九年）七月·大运河疏浚重建工程启动}

\eventanchor{ML-1411-002}{1411（永乐九年）七月6日·宝藏之旅：宝船队返回南京}

\eventanchor{ML-1411-003}{1411（永乐九年）·永乐皇帝派遣一失哈探索北满洲}

\eventanchor{ML-1411-004}{1411（永乐九年）·足利义持拒绝永乐皇帝镇压倭寇的请求}

\eventanchor{ML-1412-001}{1412（永乐十年）十二月18日·宝藏航行：永乐皇帝下令第四次下西洋}

\eventanchor{ML-1412-002}{1412（永乐十年）·明代以炮弹作为弹药。}

\eventanchor{ML-1412-003}{1412（永乐十年）·北方诸部的册封、互市或出征信息构成本年边疆关系线索，部众名称与军数须保留原材料语境。}

\eventanchor{ML-1412-004}{1412（永乐十年）·靖难后军政接收、北方防务与漕运调发的本年记录为重建政权的连续行动提供编年节点。}

\eventanchor{ML-1413-001}{1413（永乐十一年）秋·宝藏航海：郑和从南京出发，按常航路线，新增马尔代夫、比特拉、切特拉岛、霍尔木兹四个航点，并有这样的描述：“四方外国船只，并陆路来往的外国商人，都来到此地，聚集买卖，故此国民皆富”。}

\eventanchor{ML-1413-002}{1413（永乐十一年）·罗本钦波·古实日·罗卓坚赞访问南京}

\eventanchor{ML-1413-003}{1413（永乐十一年）·永宁寺碑：明朝派遣伊什哈到努尔干地区军事委员会开设驿站，传播佛教}

\eventanchor{ML-1414-001}{1414（永乐十二年）三月30日·中国第四次统治越南：陈贵圹被俘}

\eventanchor{ML-1414-002}{1414（永乐十二年）四月·第二次蒙古战役：明军在土拉河与卫拉特交战，伤亡惨重，但最终通过重炮轰击取得胜利}

\subsection{事件导航与来源边界}

\eventanchor{ML-1414-003}{1414（永乐十二年）·确杰释迦耶喜访问南京}

\eventanchor{ML-1415-001}{1415（永乐十三年）·宝藏航行：宝藏船队捕获了 Sekandar，一名反抗 Samudera Pasai 苏丹国国王 Zain al-'Abidin 的叛军}

\eventanchor{ML-1415-002}{1415（永乐十三年）六月·大运河已重建}

\eventanchor{ML-1415-003}{1415（永乐十三年）八月12日·宝藏之旅：宝船队返回南京}

\eventanchor{ML-1415-004}{1415（永乐十三年）八月13日·宝藏航行：郑和的同事被派往孟加拉运送礼物}

\eventanchor{ML-1416-001}{1416（永乐十四年）十一月19日·宝藏之旅：永乐皇帝赐礼给18国使节}

\eventanchor{ML-1416-002}{1416（永乐十四年）十二月19日·宝藏航行：永乐皇帝下令第五次下西洋}

\eventanchor{ML-1416-003}{1416（永乐十四年）·迁都筹备、官署调度和礼制安排在本年诏令与奏议中可追踪，文书形成不等于工程即时完成。}

\eventanchor{ML-1417-001}{1417（永乐十五年）·林山起义：黎来领导反抗明朝的起义}

\eventanchor{ML-1417-002}{1417（永乐十五年）秋·宝藏航行：郑和从中国出发，沿先前航线前往霍尔木兹，然后亚丁、摩加迪沙、巴拉瓦、朱布和马林迪}

\eventanchor{ML-1417-003}{1417（永乐十五年）·北方诸部的册封、互市或出征信息构成本年边疆关系线索，部众名称与军数须保留原材料语境。}

\eventanchor{ML-1418-001}{1418（永乐十六年）·靖难后军政接收、北方防务与漕运调发的本年记录为重建政权的连续行动提供编年节点。}

\eventanchor{ML-1418-002}{1418（永乐十六年）·迁都筹备、官署调度和礼制安排在本年诏令与奏议中可追踪，文书形成不等于工程即时完成。}

\eventanchor{ML-1418-003}{1418（永乐十六年）·使团、朝贡港口与航海调度的本年材料记录官方海上秩序，不能以朝贡名义替代实际贸易网络。}

\eventanchor{ML-1419-001}{1419（永乐十七年）八月8日·宝藏之旅：宝船队返回中国}

\eventanchor{ML-1419-002}{1419（永乐十七年）九月20日·宝藏之旅：大使们向明朝宫廷赠送珍奇动物，其中包括孟加拉人从索马里进口的长颈鹿}

\eventanchor{ML-1419-003}{1419（永乐十七年）十月2日·宝藏航行：下达建造41艘宝藏船的订单}

\eventanchor{ML-1419-004}{1419（永乐十七年）·元宵节期间，明皇宫会燃放烟火，火箭沿着电线运行，点燃灯笼，照亮宫殿。}

\eventanchor{ML-1420-001}{1420（永乐十八年）·紫禁城：天坛建造完成}

\eventanchor{ML-1420-002}{1420（永乐十八年）十月28日·北京正式成为明朝首都}

\eventanchor{ML-1421-001}{1421（永乐十九年）三月3日·宝藏航行：第六次下西洋的命令，向霍尔木兹等16个国家的使节赠送纸币、钱币、礼服和衬里等礼物}

\eventanchor{ML-1421-002}{1421（永乐十九年）五月14日·宝藏航行：永乐皇帝下令暂停宝藏航行}

\eventanchor{ML-1421-003}{1421（永乐十九年）十一月10日·宝藏航行：向郑和下令，提供洪堡及十六国使节回国的通道；宝藏船队按常规航线前往锡兰，在那里分头前往马尔代夫、霍尔木兹、阿拉伯国家乔法尔、拉萨和亚丁，以及两个非洲国家摩加迪沙和巴拉瓦；郑和访问甘八里}

\eventanchor{ML-1422-001}{1422（永乐二十年）·宝藏航行：宝藏船队在Samudera Pasai Sultanate重新集结并访问暹罗，然后返回中国}

\eventanchor{ML-1422-002}{1422（永乐二十年）四月·第三次蒙古战役：明朝军队出兵攻打阿鲁台，但未能与他交战并返回北京}

\eventanchor{ML-1422-003}{1422（永乐二十年）九月3日·宝藏之旅：宝船队携暹罗、巴塞苏丹国和亚丁的使节返回中国}

\eventanchor{ML-1423-001}{1423（永乐二十一年）八月·第四次蒙古战役：永乐皇帝向阿鲁台发起进攻，却发现他已经被瓦剌击败}

\eventanchor{ML-1423-002}{1423（永乐二十一年）·永乐时期典籍、学校与礼仪项目的本年记录提供知识工程线索，编纂、刻印和传播应分别核验。}

\eventanchor{ML-1424-001}{1424（永乐二十二年）二月27日·宝藏之旅：郑和出使巨港，授予史进清之子史继孙“纱帽、金纹锦袍、银印”}

\eventanchor{ML-1424-002}{1424（永乐二十二年）四月·第五次蒙古战役：永乐皇帝率军远征，讨伐阿鲁台部落的残余势力，但未能找到他们}

\eventanchor{ML-1424-003}{1424（永乐二十二年）八月12日·永乐皇帝驾崩}

\eventanchor{ML-1424-004}{1424（永乐二十二年）九月7日·宝藏之旅：朱高炽即位洪熙皇帝，宝藏之旅结束}

\eventanchor{ML-1424-005}{1424（永乐二十二年）·都市毕业生担任县长}

\eventanchor{ML-1458-003}{1458（天顺二年）·地方灾异、赈济与水利条目为本年社会史提供入口，地方志的修纂层次需要说明。}

\eventanchor{MM-1458-CROWN-PRINCE-RESTORED}{1458（天顺二年）·英宗复立朱见深为皇太子}

\eventanchor{ML-1459-001}{1459（天顺三年）·景泰、天顺时期的继承、人事与诏令链条可在本年材料中定位，政变责任不可由单一史书裁断。}

\eventanchor{ML-1459-002}{1459（天顺三年）·蒙古诸部、朝鲜及西北绿洲的交涉在本年留下多方材料，应以具体使团和地名校正概称。}

\eventanchor{ML-1459-003}{1459（天顺三年）·科举、学校和法律条文在本年制度材料中可追踪，不能据规范文本推定州县实践一致。}

\eventanchor{ML-1459-004}{1459（天顺三年）·土木危机后的边镇整顿、京师防务或复辟前后军政调度在本年材料中构成连续问题。}

\eventanchor{ML-1460-001}{1460（天顺四年）·科举、学校和法律条文在本年制度材料中可追踪，不能据规范文本推定州县实践一致。}

\eventanchor{ML-1460-002}{1460（天顺四年）·地方灾异、赈济与水利条目为本年社会史提供入口，地方志的修纂层次需要说明。}

\eventanchor{ML-1460-003}{1460（天顺四年）·景泰、天顺时期的继承、人事与诏令链条可在本年材料中定位，政变责任不可由单一史书裁断。}

\eventanchor{MM-1460-SHI-HENG-FALL}{1460（天顺四年）·石亨被逮下狱，复辟功臣集团失去中枢优势}

\eventanchor{ML-1461-001}{1461（天顺五年）八月7日·曹钦之乱：曹钦叛乱并试图攻入北京，但被捕并被迫自杀}

\eventanchor{ML-1461-002}{1461（天顺五年）·土木危机后的边镇整顿、京师防务或复辟前后军政调度在本年材料中构成连续问题。}

\eventanchor{ML-1461-003}{1461（天顺五年）·军饷、漕运和户籍赋役的本年记录揭示恢复秩序的行政成本，账面额与实收须分列。}

\eventanchor{MM-1461-CAO-QIN-REBELLION}{1461（天顺五年）七月·曹钦在北京发动兵变，攻东安门等处后被官军平定}

\eventanchor{ML-1462-001}{1462（天顺六年）·军饷、漕运和户籍赋役的本年记录揭示恢复秩序的行政成本，账面额与实收须分列。}

\eventanchor{ML-1462-002}{1462（天顺六年）·地方灾异、赈济与水利条目为本年社会史提供入口，地方志的修纂层次需要说明。}

\eventanchor{ML-1462-003}{1462（天顺六年）·景泰、天顺时期的继承、人事与诏令链条可在本年材料中定位，政变责任不可由单一史书裁断。}

\eventanchor{ML-1462-004}{1462（天顺六年）·蒙古诸部、朝鲜及西北绿洲的交涉在本年留下多方材料，应以具体使团和地名校正概称。}

\eventanchor{ML-1463-001}{1463（天顺七年）·蒙古诸部、朝鲜及西北绿洲的交涉在本年留下多方材料，应以具体使团和地名校正概称。}

\eventanchor{ML-1463-002}{1463（天顺七年）·土木危机后的边镇整顿、京师防务或复辟前后军政调度在本年材料中构成连续问题。}

\eventanchor{ML-1463-003}{1463（天顺七年）·军饷、漕运和户籍赋役的本年记录揭示恢复秩序的行政成本，账面额与实收须分列。}

\eventanchor{ML-1463-004}{1463（天顺七年）·科举、学校和法律条文在本年制度材料中可追踪，不能据规范文本推定州县实践一致。}

\eventanchor{ML-1464-001}{1464（天顺八年）二月23日·天顺帝驾崩，朱千深即位成化帝}

\eventanchor{ML-1464-002}{1464（天顺八年）·广西瑶族起义军侯大狗}

\eventanchor{ML-1464-003}{1464（天顺八年）·宝藏航行：宝藏航行的文献被刘大侠从兵部档案中移出并销毁，理由是它们是“离奇之事，远非耳目所证的欺骗性夸大”，“三宝出征西海，浪费了数十万的金钱和粮食，而且死亡的人也可以算进”。虽然他带回了许多珍贵的东西，但这对国家有什么好处呢？}

\eventanchor{ML-1506-004}{1506（正德元年）·正德皇帝乔装游走北京街头}

\eventanchor{MM-1506-EIGHT-TIGERS}{1506（正德元年）·刘瑾等八名宦官获武宗信任并进入中枢政治}

\subsection{事件导航与来源边界}

\eventanchor{ML-1507-001}{1507（正德二年）九月·元宵节花灯费35万两银子}

\eventanchor{ML-1507-002}{1507（正德二年）·科举、讲学和官员文集的本年线索可用于追踪知识网络，主张与行政结果须分开。}

\eventanchor{ML-1507-003}{1507（正德二年）·嘉靖初继统礼制的本年文书构成大礼议过程锚点，礼制理由和政治效应不可简化为单一动机。}

\eventanchor{MM-1507-LIU-JIN-FISCAL-EXTRACTIONS}{1507（正德二年）·刘瑾集团要求地方严核钱粮、差役与库藏，内廷对财政文书的介入加深}

\eventanchor{MM-1507-LIU-JIN-PURGE}{1507（正德二年）·刘瑾借厂卫和考察处置异己官员，朝廷人事与言路受到冲击}

\eventanchor{ML-1508-001}{1508（正德三年）·正德朝内廷、人事与巡幸相关文书构成本年中枢政治线索，后出道德化叙事须与原奏疏分开。（材料核查重点：诏令或奏议的形成与发受链）}

\eventanchor{ML-1508-002}{1508（正德三年）·科举、讲学和官员文集的本年线索可用于追踪知识网络，主张与行政结果须分开。}

\eventanchor{ML-1508-003}{1508（正德三年）·嘉靖初继统礼制的本年文书构成大礼议过程锚点，礼制理由和政治效应不可简化为单一动机。}

\eventanchor{ML-1508-004}{1508（正德三年）·正德朝内廷、人事与巡幸相关文书构成本年中枢政治线索，后出道德化叙事须与原奏疏分开。（材料核查重点：报称信息的来源、抄传与行政分类）}

\eventanchor{MM-1508-LIU-JIN-LAND-SURVEY}{1508（正德三年）·刘瑾推动清丈田亩和赋役整饬，以核查隐漏田土}

\eventanchor{ML-1509-001}{1509（正德四年）八月·辽东两驻军叛乱，分白银2500两后平息}

\eventanchor{ML-1509-002}{1509（正德四年）·地方反抗、边镇调兵和宗藩危机的本年军政材料标记跨区域动员，军数与战果需要复核。}

\eventanchor{ML-1509-003}{1509（正德四年）·嘉靖初继统礼制的本年文书构成大礼议过程锚点，礼制理由和政治效应不可简化为单一动机。}

\eventanchor{ML-1509-004}{1509（正德四年）·正德朝内廷、人事与巡幸相关文书构成本年中枢政治线索，后出道德化叙事须与原奏疏分开。}

\eventanchor{ML-1510-001}{1510（正德五年）五月12日·安化王叛乱：山西朱之璠叛乱}

\eventanchor{ML-1510-002}{1510（正德五年）五月30日·安化王叛乱：朱之璠被俘}

\eventanchor{ML-1510-003}{1510（正德五年）·达延汗征服鄂尔多斯河套地区}

\eventanchor{MM-1510-ANHUA-REBELLION}{1510（正德五年）五月·安化王朱寘鐇在宁夏起兵，旋被当地官军与将领平定}

\eventanchor{MM-1510-LIU-JIN-EXECUTION}{1510（正德五年）八月·刘瑾被以谋反等罪处死}

\eventanchor{ML-1511-001}{1511（正德六年）二月·北京周边土匪起义}

\eventanchor{ML-1511-002}{1511（正德六年）十月·北京各地土匪烧毁皇家粮车}

\eventanchor{ML-1511-003}{1511（正德六年）·占领马六甲 (1511)：马六甲苏丹国向葡萄牙人发出求援请求}

\eventanchor{MM-1511-LIU-LIU-QI-UPRISING}{1511（正德六年）·刘六、刘七等在山东、北直隶一带聚众活动，成为跨区域治安危机}

\eventanchor{ML-1512-001}{1512（正德七年）一月·土匪袭击霸州}

\eventanchor{ML-1512-002}{1512（正德七年）九月7日·强盗大军被击败}

\eventanchor{ML-1512-003}{1512（正德七年）·地方反抗、边镇调兵和宗藩危机的本年军政材料标记跨区域动员，军数与战果需要复核。}

\eventanchor{MM-1512-LIU-LIU-QI-SUPPRESSION}{1512（正德七年）·官军在山东、北直隶分路追击刘六、刘七等集团，主要武装至年内瓦解}

\eventanchor{ML-1513-001}{1513（正德八年）·清丈、库藏、差役与征解的本年记录显示财政监管压力，整饬命令不等于隐漏已被清除。（材料核查重点：诏令或奏议的形成与发受链）}

\eventanchor{ML-1513-002}{1513（正德八年）·正德朝内廷、人事与巡幸相关文书构成本年中枢政治线索，后出道德化叙事须与原奏疏分开。}

\eventanchor{ML-1513-003}{1513（正德八年）·清丈、库藏、差役与征解的本年记录显示财政监管压力，整饬命令不等于隐漏已被清除。（材料核查重点：报称信息的来源、抄传与行政分类）}

\eventanchor{MM-1513-PORTUGUESE-TAMAO-CONTACT}{1513（正德八年）·葡萄牙航海者在屯门一带停泊并尝试贸易，为后续使团接触提供海上先例}

\eventanchor{ML-1514-001}{1514（正德九年）二月10日·宫殿庭院内的火药帐篷着火，烧毁了住宅宫殿}

\eventanchor{ML-1514-002}{1514（正德九年）九月·正德皇帝被老虎咬成重伤}

\eventanchor{ML-1514-003}{1514（正德九年）·东南港口、日本使团或葡萄牙接触的本年材料为海上交往留档，朝贡、私贸与冲突不可混称。}

\eventanchor{MM-1514-TURPAN-HAMI-ANNEXATION}{1514（正德九年）·吐鲁番势力再次夺取哈密，明朝维系该绿洲朝贡政权的能力显著下降}

\eventanchor{ML-1515-001}{1515（正德十年）夏·京师、侍卫三万出动，重建宫殿}

\eventanchor{ML-1515-002}{1515（正德十年）·东南港口、日本使团或葡萄牙接触的本年材料为海上交往留档，朝贡、私贸与冲突不可混称。}

\eventanchor{ML-1515-003}{1515（正德十年）·地方反抗、边镇调兵和宗藩危机的本年军政材料标记跨区域动员，军数与战果需要复核。}

\eventanchor{ML-1515-004}{1515（正德十年）·科举、讲学和官员文集的本年线索可用于追踪知识网络，主张与行政结果须分开。}

\eventanchor{ML-1516-001}{1516（正德十一年）·地方反抗、边镇调兵和宗藩危机的本年军政材料标记跨区域动员，军数与战果需要复核。}

\eventanchor{ML-1516-002}{1516（正德十一年）·科举、讲学和官员文集的本年线索可用于追踪知识网络，主张与行政结果须分开。}

\eventanchor{ML-1516-003}{1516（正德十一年）·清丈、库藏、差役与征解的本年记录显示财政监管压力，整饬命令不等于隐漏已被清除。}

\eventanchor{ML-1516-004}{1516（正德十一年）·嘉靖初继统礼制的本年文书构成大礼议过程锚点，礼制理由和政治效应不可简化为单一动机。}

\eventanchor{ML-1517-001}{1517（正德十二年）十月16日·达延汗袭击明朝}

\eventanchor{ML-1517-002}{1517（正德十二年）十月20日·正德皇帝击退达延汗的袭击}

\eventanchor{ML-1517-003}{1517（正德十二年）·托梅·皮雷抵达广州}

\eventanchor{MM-1517-PORTUGUESE-EMBASSY}{1517（正德十二年）·葡萄牙使团抵达广东，托名满剌加使者请求入朝}

\eventanchor{ML-1518-001}{1518（正德十三年）一月·正德皇帝因没有给他足够的钱财而被囚禁在北京}

\eventanchor{ML-1518-002}{1518（正德十三年）·东南港口、日本使团或葡萄牙接触的本年材料为海上交往留档，朝贡、私贸与冲突不可混称。}

\eventanchor{ML-1518-003}{1518（正德十三年）·正德朝内廷、人事与巡幸相关文书构成本年中枢政治线索，后出道德化叙事须与原奏疏分开。}

\eventanchor{MM-1518-WANG-YANGMING-GANZHOU}{1518（正德十三年）·王守仁在赣南主持军政，推行保甲、招抚与军事征剿并用}

\eventanchor{ML-1519-001}{1519（正德十四年）七月9日·宁王叛乱：江西朱辰灏叛乱}

\eventanchor{ML-1519-002}{1519（正德十四年）七月13日·宁王叛乱：叛军攻克九江}

\eventanchor{ML-1519-003}{1519（正德十四年）七月23日·宁王叛乱：叛军围攻安庆}

\eventanchor{ML-1519-004}{1519（正德十四年）八月9日·宁王叛乱：叛军解围安庆}

\eventanchor{ML-1519-005}{1519（正德十四年）八月13日·宁王叛乱：清军攻克南昌}

\eventanchor{ML-1519-006}{1519（正德十四年）八月15日·宁王叛乱：朱辰灏军大败}

\eventanchor{ML-1519-007}{1519（正德十四年）八月20日·宁王叛乱：朱辰浩逃离舰队被俘}

\eventanchor{MM-1519-NING-REBELLION}{1519（正德十四年）六月·宁王朱宸濠在南昌起兵，试图沿江扩张}

\eventanchor{MM-1519-WANG-YANGMING-VICTORY}{1519（正德十四年）七月·王守仁率地方军在鄱阳湖、南昌一带击败并俘获朱宸濠}

\subsection{事件导航与来源边界}

\eventanchor{MM-1519-WUZONG-SOUTHERN-TOUR}{1519（正德十四年）·宁王之变已受控制后，武宗仍南巡并以亲临军务为名驻跸江南}

\eventanchor{ML-1520-001}{1520（正德十五年）一月·正德皇帝禁止宰猪}

\eventanchor{ML-1520-002}{1520（正德十五年）五月·葡萄牙人贿赂了广州的一名太监，让他们通过，托梅·皮雷一行抵达南京}

\eventanchor{ML-1520-003}{1520（正德十五年）·科举、讲学和官员文集的本年线索可用于追踪知识网络，主张与行政结果须分开。}

\eventanchor{ML-1520-004}{1520（正德十五年）·地方反抗、边镇调兵和宗藩危机的本年军政材料标记跨区域动员，军数与战果需要复核。}

\eventanchor{ML-1521-001}{1521（正德十六年）四月20日·正德皇帝驾崩}

\eventanchor{ML-1521-002}{1521（正德十六年）四月21日·托梅·皮雷的政党被驱逐出北京}

\eventanchor{ML-1521-003}{1521（正德十六年）四月27日·朱厚熜即位嘉靖皇帝}

\eventanchor{ML-1521-004}{1521（正德十六年）五月·屯门之战：明军将拒绝离开的葡萄牙舰队驱逐出屯门}

\eventanchor{ML-1521-005}{1521（正德十六年）·宫殿重建完成}

\eventanchor{MM-1521-TUNMEN-CONFLICT}{1521（正德十六年）·屯门海域发生明军与葡萄牙船队冲突}

\eventanchor{ML-1558-004}{1558（嘉靖三十七年）·帝国国库白银跌至不足20万盎司}

\eventanchor{MM-1558-QI-JIGUANG-TRAINING}{1558（嘉靖三十七年）·戚继光在浙江募练义乌等地兵员，形成适应沿海作战的部队}

\eventanchor{ML-1559-001}{1559（嘉靖三十八年）夏·戚继光开始对新兵进行战术改革}

\eventanchor{ML-1559-002}{1559（嘉靖三十八年）·长江三角洲地区干旱导致饥荒}

\eventanchor{ML-1559-003}{1559（嘉靖三十八年）十二月·苏州卫戍兵变}

\eventanchor{ML-1559-004}{1559（嘉靖三十八年）·嘉靖倭口劫掠：海盗占领金门岛并袭击福建和广东}

\eventanchor{MM-1559-HU-ZONGXIAN-SEA-DEFENSE}{1559（嘉靖三十八年）·胡宗宪督理浙直军务，以招抚、离间和军事追击并用应对海上武装}

\eventanchor{MM-1559-WANG-ZHI-EXECUTION}{1559（嘉靖三十八年）·王直在杭州被处决，明廷以司法处置结束其招抚案件}

\eventanchor{ML-1560-001}{1560（嘉靖三十九年）三月·南京卫戍叛军回应削减口粮，直到给他们4万两白银}

\eventanchor{ML-1560-002}{1560（嘉靖三十九年）·嘉靖皇帝患有失眠症}

\eventanchor{ML-1560-003}{1560（嘉靖三十九年）·戚继光发表了《急效新书》，描述了火枪齐射技术以及他训练明军使用该技术的经验。}

\eventanchor{ML-1560-004}{1560（嘉靖三十九年）·封贡、互市、月港和港口管理的本年文书记录边海关系重组，官方许可不等于贸易全貌。}

\eventanchor{ML-1560-005}{1560（嘉靖三十九年）·严嵩时期至隆庆初的人事和奏议材料标记中枢权力变化，案狱与党争标签需保留来源立场。}

\eventanchor{ML-1561-001}{1561（嘉靖四十年）十二月·紫禁城的住宅宫殿在一场大火中被毁}

\eventanchor{ML-1561-002}{1561（嘉靖四十年）·明朝开始生产便携式后装火器。}

\eventanchor{ML-1561-003}{1561（嘉靖四十年）·严嵩时期至隆庆初的人事和奏议材料标记中枢权力变化，案狱与党争标签需保留来源立场。}

\eventanchor{ML-1561-004}{1561（嘉靖四十年）·海防知识、学校、科举与礼制的本年文本提供制度和知识史线索，方案不等于实际配置。}

\eventanchor{MM-1561-TAIZHOU-VICTORIES}{1561（嘉靖四十年）·戚继光部在台州一线连续击败海上武装，巩固浙东防务}

\eventanchor{ML-1562-001}{1562（嘉靖四十一年）六月·故宫的住宅宫殿被重建}

\eventanchor{ML-1562-002}{1562（嘉靖四十一年）十二月·嘉靖倭口劫掠：海盗占领兴化}

\eventanchor{ML-1562-003}{1562（嘉靖四十一年）·海防知识、学校、科举与礼制的本年文本提供制度和知识史线索，方案不等于实际配置。}

\eventanchor{ML-1562-004}{1562（嘉靖四十一年）·本年灾情、赈恤或河工记录连接中央与地方社会，区域差异需以府县材料追踪。}

\eventanchor{MM-1562-FUJIAN-CAMPAIGNS}{1562（嘉靖四十一年）·戚继光等率军进入福建，围攻并清除若干沿海武装据点}

\eventanchor{ML-1563-001}{1563（嘉靖四十二年）五月·嘉靖倭寇突袭：明军夺回兴化并摧毁福建海盗基地}

\eventanchor{ML-1563-002}{1563（嘉靖四十二年）·海盗林道前遭明军追击，撤退至台湾西南部}

\eventanchor{ML-1563-003}{1563（嘉靖四十二年）·明朝将军下令在澎湖修建城墙}

\eventanchor{ML-1563-004}{1563（嘉靖四十二年）·俺答边患、东南海防或沿海武装的本年军令与战报构成危机处理线索，战果和参与者须交叉核对。}

\eventanchor{MM-1563-PINGHAI-RECOVERY}{1563（嘉靖四十二年）·官军收复平海卫等福建沿海据点}

\eventanchor{ML-1564-001}{1564（嘉靖四十三年）·太监把桃子扔到嘉靖皇帝的床上并告诉他这是从天上掉下来的}

\eventanchor{ML-1564-002}{1564（嘉靖四十三年）·嘉靖皇帝因应俸禄要求，将所有宗族降为平民}

\eventanchor{ML-1564-003}{1564（嘉靖四十三年）·本年灾情、赈恤或河工记录连接中央与地方社会，区域差异需以府县材料追踪。}

\eventanchor{ML-1564-004}{1564（嘉靖四十三年）·军饷、盐课、漕运和赋役的本年行政材料显示中晚明财政压力，法定额与实收不可互代。}

\eventanchor{MM-1564-QI-JIGUANG-FUJIAN-COMMAND}{1564（嘉靖四十三年）·戚继光继续统领福建陆路军务，与俞大猷等分担沿海据点的清剿和守备}

\eventanchor{ML-1565-001}{1565（嘉靖四十四年）·嘉靖皇帝病重}

\eventanchor{ML-1565-002}{1565（嘉靖四十四年）·本年灾情、赈恤或河工记录连接中央与地方社会，区域差异需以府县材料追踪。}

\eventanchor{ML-1565-003}{1565（嘉靖四十四年）·俺答边患、东南海防或沿海武装的本年军令与战报构成危机处理线索，战果和参与者须交叉核对。}

\eventanchor{ML-1565-004}{1565（嘉靖四十四年）·封贡、互市、月港和港口管理的本年文书记录边海关系重组，官方许可不等于贸易全貌。}

\eventanchor{MM-1565-HU-ZONGXIAN-ARREST}{1565（嘉靖四十四年）·严嵩失势后，胡宗宪被逮治并死于狱中}

\eventanchor{ML-1566-001}{1566（嘉靖四十五年）·嘉靖倭寇劫掠：明军铲除江西、广东海盗}

\eventanchor{ML-1566-002}{1566（嘉靖四十五年）·俺答边患、东南海防或沿海武装的本年军令与战报构成危机处理线索，战果和参与者须交叉核对。}

\eventanchor{ML-1566-003}{1566（嘉靖四十五年）·军饷、盐课、漕运和赋役的本年行政材料显示中晚明财政压力，法定额与实收不可互代。}

\eventanchor{ML-1566-004}{1566（嘉靖四十五年）·严嵩时期至隆庆初的人事和奏议材料标记中枢权力变化，案狱与党争标签需保留来源立场。}

\eventanchor{MM-1566-HAI-RUI-MEMORIAL}{1566（嘉靖四十五年）·海瑞上《治安疏》批评嘉靖政治，随后被下狱}

\eventanchor{ML-1567-001}{1567（隆庆元年）一月23日·嘉靖皇帝驾崩}

\eventanchor{ML-1567-002}{1567（隆庆元年）二月4日·朱载厚即位隆庆皇帝}

\eventanchor{ML-1567-003}{1567（隆庆元年）·海外交易禁令解除}

\eventanchor{MM-1567-HAI-RUI-RELEASE}{1567（隆庆元年）正月·隆庆新朝释放因治安疏入狱的海瑞，撤销嘉靖末对其的拘系}

\eventanchor{MM-1567-LONGQING-ACCESSION}{1567（隆庆元年）正月·嘉靖帝死后，皇太子朱载垕即位，是为隆庆帝}

\eventanchor{MM-1567-YUEGANG-OPENING}{1567（隆庆元年）·福建月港获准在限定条件下办理民间海外贸易}

\subsection{事件导航与来源边界}

\eventanchor{MM-1567-ZHANG-JUZHENG-ENTERED-GRAND-SECRETARIAT}{1567（隆庆元年）·张居正入阁，参与隆庆朝边务、财政与人事协调}

\eventanchor{ML-1568-001}{1568（隆庆二年）·俺答边患、东南海防或沿海武装的本年军令与战报构成危机处理线索，战果和参与者须交叉核对。}

\eventanchor{ML-1568-002}{1568（隆庆二年）·军饷、盐课、漕运和赋役的本年行政材料显示中晚明财政压力，法定额与实收不可互代。}

\eventanchor{ML-1568-003}{1568（隆庆二年）·封贡、互市、月港和港口管理的本年文书记录边海关系重组，官方许可不等于贸易全貌。}

\eventanchor{ML-1568-004}{1568（隆庆二年）·海防知识、学校、科举与礼制的本年文本提供制度和知识史线索，方案不等于实际配置。}

\eventanchor{MM-1568-QI-JIGUANG-JIZHOU}{1568（隆庆二年）·戚继光调任蓟州总兵，开始整饬蓟镇练兵与防务}

\eventanchor{ML-1569-001}{1569（隆庆三年）·军饷、盐课、漕运和赋役的本年行政材料显示中晚明财政压力，法定额与实收不可互代。}

\eventanchor{ML-1569-002}{1569（隆庆三年）·封贡、互市、月港和港口管理的本年文书记录边海关系重组，官方许可不等于贸易全貌。}

\eventanchor{ML-1569-003}{1569（隆庆三年）·严嵩时期至隆庆初的人事和奏议材料标记中枢权力变化，案狱与党争标签需保留来源立场。}

\eventanchor{ML-1569-004}{1569（隆庆三年）·本年灾情、赈恤或河工记录连接中央与地方社会，区域差异需以府县材料追踪。}

\eventanchor{MM-1569-HAI-RUI-SOUTHERN-METROPOLIS}{1569（隆庆三年）·海瑞出任巡抚应天十府，处理积案、赋役与土地争讼}

\eventanchor{ML-1570-001}{1570（隆庆四年）·建州卫王高袭击明朝聚居地}

\eventanchor{ML-1570-002}{1570（隆庆四年）·严嵩时期至隆庆初的人事和奏议材料标记中枢权力变化，案狱与党争标签需保留来源立场。}

\eventanchor{ML-1570-003}{1570（隆庆四年）·海防知识、学校、科举与礼制的本年文本提供制度和知识史线索，方案不等于实际配置。}

\eventanchor{ML-1570-004}{1570（隆庆四年）·俺答边患、东南海防或沿海武装的本年军令与战报构成危机处理线索，战果和参与者须交叉核对。}

\eventanchor{MM-1570-ALTAN-NEGOTIATIONS}{1570（隆庆四年）·明廷与俺答部围绕封贡互市及赵全等人问题展开谈判}

\eventanchor{ML-1571-001}{1571（隆庆五年）·封贡、互市、月港和港口管理的本年文书记录边海关系重组，官方许可不等于贸易全貌。}

\eventanchor{ML-1571-002}{1571（隆庆五年）·海防知识、学校、科举与礼制的本年文本提供制度和知识史线索，方案不等于实际配置。}

\eventanchor{ML-1571-003}{1571（隆庆五年）·本年灾情、赈恤或河工记录连接中央与地方社会，区域差异需以府县材料追踪。}

\eventanchor{ML-1571-004}{1571（隆庆五年）·军饷、盐课、漕运和赋役的本年行政材料显示中晚明财政压力，法定额与实收不可互代。}

\eventanchor{MM-1571-LONGQING-PEACE}{1571（隆庆五年）·明廷封俺答为顺义王，并在宣大等处恢复互市}

\eventanchor{ML-1572-001}{1572（隆庆六年）七月5日·隆庆皇帝驾崩}

\eventanchor{ML-1572-002}{1572（隆庆六年）七月19日·朱翊钧成为万历皇帝}

\eventanchor{ML-1572-003}{1572（隆庆六年）·封贡、互市、月港和港口管理的本年文书记录边海关系重组，官方许可不等于贸易全貌。}

\eventanchor{ML-1606-002}{1606（万历三十四年）·明式步枪均附有插入式刺刀。}

\eventanchor{ML-1606-003}{1606（万历三十四年）·后金兴起、朝鲜交涉和海上网络的本年多方材料揭示区域关系变化，政权名称与译名须按原文说明。（材料核查重点：诏令或奏议的形成与发受链）}

\eventanchor{ML-1606-004}{1606（万历三十四年）·本年灾害、疫病或地方赈济线索应以受灾地区文书和地方志复核，中央奏报不能覆盖全部社会。}

\eventanchor{ML-1606-005}{1606（万历三十四年）·后金兴起、朝鲜交涉和海上网络的本年多方材料揭示区域关系变化，政权名称与译名须按原文说明。（材料核查重点：报称信息的来源、抄传与行政分类）}

\eventanchor{ML-1606-006}{1606（万历三十四年）·辽东战事、边镇防务和西南兵事的本年军令记录资源调动，战报数字应与朝鲜、满文或地方材料比照。}

\eventanchor{ML-1607-001}{1607（万历三十五年）·欧几里得几何原理前六册被翻译成中文}

\eventanchor{ML-1607-002}{1607（万历三十五年）·本年灾害、疫病或地方赈济线索应以受灾地区文书和地方志复核，中央奏报不能覆盖全部社会。}

\eventanchor{ML-1607-003}{1607（万历三十五年）·东林、书院、刻书和宗教活动的本年材料记录精英网络，社团存在不等于一致政治立场。（材料核查重点：诏令或奏议的形成与发受链）}

\eventanchor{ML-1607-004}{1607（万历三十五年）·万历末至天启初的立储、内阁和言路材料标记中枢运作的堵塞与调整，派系归类需回到具体文书。}

\eventanchor{ML-1607-005}{1607（万历三十五年）·东林、书院、刻书和宗教活动的本年材料记录精英网络，社团存在不等于一致政治立场。（材料核查重点：报称信息的来源、抄传与行政分类）}

\eventanchor{ML-1607-006}{1607（万历三十五年）·后金兴起、朝鲜交涉和海上网络的本年多方材料揭示区域关系变化，政权名称与译名须按原文说明。}

\eventanchor{ML-1608-001}{1608（万历三十六年）·矿税、加派、漕运和辽饷的本年奏疏与账目提供财政压力线索，预算、欠额和实解不得混同。（材料核查重点：诏令或奏议的形成与发受链）}

\eventanchor{ML-1608-002}{1608（万历三十六年）·万历末至天启初的立储、内阁和言路材料标记中枢运作的堵塞与调整，派系归类需回到具体文书。}

\eventanchor{ML-1608-003}{1608（万历三十六年）·本年灾害、疫病或地方赈济线索应以受灾地区文书和地方志复核，中央奏报不能覆盖全部社会。（材料核查重点：诏令或奏议的形成与发受链）}

\eventanchor{ML-1608-004}{1608（万历三十六年）·矿税、加派、漕运和辽饷的本年奏疏与账目提供财政压力线索，预算、欠额和实解不得混同。（材料核查重点：报称信息的来源、抄传与行政分类）}

\eventanchor{ML-1608-005}{1608（万历三十六年）·本年灾害、疫病或地方赈济线索应以受灾地区文书和地方志复核，中央奏报不能覆盖全部社会。（材料核查重点：报称信息的来源、抄传与行政分类）}

\eventanchor{ML-1608-006}{1608（万历三十六年）·东林、书院、刻书和宗教活动的本年材料记录精英网络，社团存在不等于一致政治立场。}

\eventanchor{ML-1609-001}{1609（万历三十七年）·贾运河竣工}

\eventanchor{ML-1609-002}{1609（万历三十七年）·矿税、加派、漕运和辽饷的本年奏疏与账目提供财政压力线索，预算、欠额和实解不得混同。}

\eventanchor{ML-1609-003}{1609（万历三十七年）·万历末至天启初的立储、内阁和言路材料标记中枢运作的堵塞与调整，派系归类需回到具体文书。（材料核查重点：诏令或奏议的形成与发受链）}

\eventanchor{ML-1609-004}{1609（万历三十七年）·辽东战事、边镇防务和西南兵事的本年军令记录资源调动，战报数字应与朝鲜、满文或地方材料比照。}

\eventanchor{ML-1609-005}{1609（万历三十七年）·万历末至天启初的立储、内阁和言路材料标记中枢运作的堵塞与调整，派系归类需回到具体文书。（材料核查重点：报称信息的来源、抄传与行政分类）}

\eventanchor{ML-1609-006}{1609（万历三十七年）·本年灾害、疫病或地方赈济线索应以受灾地区文书和地方志复核，中央奏报不能覆盖全部社会。}

\eventanchor{ML-1610-001}{1610（万历三十八年）·李约瑟估计，欧洲文明在天文学和物理学方面大约在这个时候超越了中国}

\eventanchor{ML-1610-002}{1610（万历三十八年）·辽东战事、边镇防务和西南兵事的本年军令记录资源调动，战报数字应与朝鲜、满文或地方材料比照。}

\eventanchor{ML-1610-003}{1610（万历三十八年）·矿税、加派、漕运和辽饷的本年奏疏与账目提供财政压力线索，预算、欠额和实解不得混同。（材料核查重点：诏令或奏议的形成与发受链）}

\eventanchor{ML-1610-004}{1610（万历三十八年）·后金兴起、朝鲜交涉和海上网络的本年多方材料揭示区域关系变化，政权名称与译名须按原文说明。}

\eventanchor{ML-1610-005}{1610（万历三十八年）·矿税、加派、漕运和辽饷的本年奏疏与账目提供财政压力线索，预算、欠额和实解不得混同。（材料核查重点：报称信息的来源、抄传与行政分类）}

\eventanchor{ML-1610-006}{1610（万历三十八年）·万历末至天启初的立储、内阁和言路材料标记中枢运作的堵塞与调整，派系归类需回到具体文书。}

\eventanchor{ML-1611-001}{1611（万历三十九年）·辽东战事、边镇防务和西南兵事的本年军令记录资源调动，战报数字应与朝鲜、满文或地方材料比照。（材料核查重点：诏令或奏议的形成与发受链）}

\eventanchor{ML-1611-002}{1611（万历三十九年）·后金兴起、朝鲜交涉和海上网络的本年多方材料揭示区域关系变化，政权名称与译名须按原文说明。}

\eventanchor{ML-1611-003}{1611（万历三十九年）·辽东战事、边镇防务和西南兵事的本年军令记录资源调动，战报数字应与朝鲜、满文或地方材料比照。（材料核查重点：报称信息的来源、抄传与行政分类）}

\eventanchor{ML-1611-004}{1611（万历三十九年）·东林、书院、刻书和宗教活动的本年材料记录精英网络，社团存在不等于一致政治立场。}

\eventanchor{ML-1611-005}{1611（万历三十九年）·辽东战事、边镇防务和西南兵事的本年军令记录资源调动，战报数字应与朝鲜、满文或地方材料比照。（材料核查重点：同年地方材料所见的执行范围）}

\eventanchor{ML-1611-006}{1611（万历三十九年）·矿税、加派、漕运和辽饷的本年奏疏与账目提供财政压力线索，预算、欠额和实解不得混同。}

\eventanchor{ML-1612-001}{1612（万历四十年）·后金兴起、朝鲜交涉和海上网络的本年多方材料揭示区域关系变化，政权名称与译名须按原文说明。（材料核查重点：诏令或奏议的形成与发受链）}

\subsection{事件导航与来源边界}

\eventanchor{ML-1612-002}{1612（万历四十年）·东林、书院、刻书和宗教活动的本年材料记录精英网络，社团存在不等于一致政治立场。}

\eventanchor{ML-1612-003}{1612（万历四十年）·后金兴起、朝鲜交涉和海上网络的本年多方材料揭示区域关系变化，政权名称与译名须按原文说明。（材料核查重点：报称信息的来源、抄传与行政分类）}

\eventanchor{ML-1612-004}{1612（万历四十年）·本年灾害、疫病或地方赈济线索应以受灾地区文书和地方志复核，中央奏报不能覆盖全部社会。}

\eventanchor{ML-1612-005}{1612（万历四十年）·后金兴起、朝鲜交涉和海上网络的本年多方材料揭示区域关系变化，政权名称与译名须按原文说明。（材料核查重点：同年地方材料所见的执行范围）}

\eventanchor{ML-1612-006}{1612（万历四十年）·辽东战事、边镇防务和西南兵事的本年军令记录资源调动，战报数字应与朝鲜、满文或地方材料比照。}

\eventanchor{ML-1613-001}{1613（万历四十一年）·东林、书院、刻书和宗教活动的本年材料记录精英网络，社团存在不等于一致政治立场。（材料核查重点：诏令或奏议的形成与发受链）}

\eventanchor{ML-1613-002}{1613（万历四十一年）·本年灾害、疫病或地方赈济线索应以受灾地区文书和地方志复核，中央奏报不能覆盖全部社会。}

\eventanchor{ML-1613-003}{1613（万历四十一年）·东林、书院、刻书和宗教活动的本年材料记录精英网络，社团存在不等于一致政治立场。（材料核查重点：报称信息的来源、抄传与行政分类）}

\eventanchor{ML-1613-004}{1613（万历四十一年）·万历末至天启初的立储、内阁和言路材料标记中枢运作的堵塞与调整，派系归类需回到具体文书。}

\eventanchor{ML-1613-005}{1613（万历四十一年）·东林、书院、刻书和宗教活动的本年材料记录精英网络，社团存在不等于一致政治立场。（材料核查重点：同年地方材料所见的执行范围）}

\eventanchor{ML-1613-006}{1613（万历四十一年）·后金兴起、朝鲜交涉和海上网络的本年多方材料揭示区域关系变化，政权名称与译名须按原文说明。}

\eventanchor{ML-1614-001}{1614（万历四十二年）·本年灾害、疫病或地方赈济线索应以受灾地区文书和地方志复核，中央奏报不能覆盖全部社会。（材料核查重点：诏令或奏议的形成与发受链）}

\eventanchor{ML-1614-002}{1614（万历四十二年）·万历末至天启初的立储、内阁和言路材料标记中枢运作的堵塞与调整，派系归类需回到具体文书。}

\eventanchor{ML-1614-003}{1614（万历四十二年）·本年灾害、疫病或地方赈济线索应以受灾地区文书和地方志复核，中央奏报不能覆盖全部社会。（材料核查重点：报称信息的来源、抄传与行政分类）}

\eventanchor{ML-1614-004}{1614（万历四十二年）·矿税、加派、漕运和辽饷的本年奏疏与账目提供财政压力线索，预算、欠额和实解不得混同。}

\eventanchor{ML-1614-005}{1614（万历四十二年）·本年灾害、疫病或地方赈济线索应以受灾地区文书和地方志复核，中央奏报不能覆盖全部社会。（材料核查重点：同年地方材料所见的执行范围）}

\eventanchor{ML-1614-006}{1614（万历四十二年）·东林、书院、刻书和宗教活动的本年材料记录精英网络，社团存在不等于一致政治立场。}

\eventanchor{ML-1615-001}{1615（万历四十三年）·努尔哈赤派遣最后一位朝贡使者前往北京}

\eventanchor{ML-1615-002}{1615（万历四十三年）·矿税、加派、漕运和辽饷的本年奏疏与账目提供财政压力线索，预算、欠额和实解不得混同。}

\eventanchor{ML-1615-003}{1615（万历四十三年）·万历末至天启初的立储、内阁和言路材料标记中枢运作的堵塞与调整，派系归类需回到具体文书。（材料核查重点：诏令或奏议的形成与发受链）}

\eventanchor{ML-1615-004}{1615（万历四十三年）·辽东战事、边镇防务和西南兵事的本年军令记录资源调动，战报数字应与朝鲜、满文或地方材料比照。}

\eventanchor{ML-1615-005}{1615（万历四十三年）·万历末至天启初的立储、内阁和言路材料标记中枢运作的堵塞与调整，派系归类需回到具体文书。（材料核查重点：报称信息的来源、抄传与行政分类）}

\eventanchor{ML-1615-006}{1615（万历四十三年）·本年灾害、疫病或地方赈济线索应以受灾地区文书和地方志复核，中央奏报不能覆盖全部社会。}

\eventanchor{ML-1616-001}{1616（万历四十四年）·努尔哈赤宣告后金，又称阿玛嘎爱新古伦}

\eventanchor{ML-1616-002}{1616（万历四十四年）·辽东战事、边镇防务和西南兵事的本年军令记录资源调动，战报数字应与朝鲜、满文或地方材料比照。}

\eventanchor{ML-1616-003}{1616（万历四十四年）·矿税、加派、漕运和辽饷的本年奏疏与账目提供财政压力线索，预算、欠额和实解不得混同。（材料核查重点：诏令或奏议的形成与发受链）}

\eventanchor{ML-1616-004}{1616（万历四十四年）·后金兴起、朝鲜交涉和海上网络的本年多方材料揭示区域关系变化，政权名称与译名须按原文说明。}

\eventanchor{ML-1616-005}{1616（万历四十四年）·矿税、加派、漕运和辽饷的本年奏疏与账目提供财政压力线索，预算、欠额和实解不得混同。（材料核查重点：报称信息的来源、抄传与行政分类）}

\eventanchor{ML-1616-006}{1616（万历四十四年）·万历末至天启初的立储、内阁和言路材料标记中枢运作的堵塞与调整，派系归类需回到具体文书。}

\eventanchor{ML-1617-001}{1617（万历四十五年）·万历末至天启初的立储、内阁和言路材料标记中枢运作的堵塞与调整，派系归类需回到具体文书。}

\eventanchor{ML-1617-002}{1617（万历四十五年）·后金兴起、朝鲜交涉和海上网络的本年多方材料揭示区域关系变化，政权名称与译名须按原文说明。}

\eventanchor{ML-1617-003}{1617（万历四十五年）·辽东战事、边镇防务和西南兵事的本年军令记录资源调动，战报数字应与朝鲜、满文或地方材料比照。（材料核查重点：诏令或奏议的形成与发受链）}

\eventanchor{ML-1617-004}{1617（万历四十五年）·东林、书院、刻书和宗教活动的本年材料记录精英网络，社团存在不等于一致政治立场。}

\eventanchor{ML-1617-005}{1617（万历四十五年）·辽东战事、边镇防务和西南兵事的本年军令记录资源调动，战报数字应与朝鲜、满文或地方材料比照。（材料核查重点：报称信息的来源、抄传与行政分类）}

\eventanchor{ML-1617-006}{1617（万历四十五年）·矿税、加派、漕运和辽饷的本年奏疏与账目提供财政压力线索，预算、欠额和实解不得混同。}

\eventanchor{ML-1618-001}{1618（万历四十六年）五月9日·抚顺之战：后金占领抚顺}

\eventanchor{ML-1618-002}{1618（万历四十六年）夏·清河之战：后金攻克清河}

\eventanchor{ML-1618-003}{1618（万历四十六年）·后金兴起、朝鲜交涉和海上网络的本年多方材料揭示区域关系变化，政权名称与译名须按原文说明。（材料核查重点：诏令或奏议的形成与发受链）}

\eventanchor{ML-1618-004}{1618（万历四十六年）·本年灾害、疫病或地方赈济线索应以受灾地区文书和地方志复核，中央奏报不能覆盖全部社会。}

\eventanchor{ML-1618-005}{1618（万历四十六年）·后金兴起、朝鲜交涉和海上网络的本年多方材料揭示区域关系变化，政权名称与译名须按原文说明。（材料核查重点：报称信息的来源、抄传与行政分类）}

\eventanchor{ML-1619-001}{1619（万历四十七年）四月18日·萨尔浒之战：明军被后金歼灭}

\eventanchor{ML-1619-002}{1619（万历四十七年）七月26日·开原之战：后金攻占开原}

\eventanchor{ML-1619-003}{1619（万历四十七年）九月3日·铁岭之战：后金占领铁岭}

\eventanchor{ML-1619-004}{1619（万历四十七年）九月·首位俄罗斯特使伊万·佩特林抵达北京}

\eventanchor{ML-1619-005}{1619（万历四十七年）·东林、书院、刻书和宗教活动的本年材料记录精英网络，社团存在不等于一致政治立场。}

\eventanchor{ML-1620-001}{1620（万历四十八年）八月18日·万历皇帝驾崩}

\eventanchor{ML-1620-002}{1620（万历四十八年）八月28日·朱常洛即位太常皇帝}

\eventanchor{ML-1620-003}{1620（万历四十八年）九月6日·太常皇帝病了}

\eventanchor{ML-1620-004}{1620（万历四十八年）九月26日·太常帝驾崩}

\eventanchor{ML-1620-005}{1620（万历四十八年）十月1日·朱由娇成为天启皇帝}

\eventanchor{ML-1620-006}{1620（万历四十八年）·明代铸造厂开始生产红衣袍。}

\eventanchor{ML-1640-006}{1640（崇祯十三年）·辽饷、剿饷、练饷及地方征解的本年文书显示财政负担，定额、征收与实际流向不得混为一谈。（材料核查重点：同年地方材料所见的执行范围）}

\eventanchor{ML-1640-007}{1640（崇祯十三年）·朝鲜、辽东和海上关系的本年材料提供外部观察位置，明、后金（清）与地方势力的称谓需具体化。}

\eventanchor{ML-1641-001}{1641（崇祯十四年）·李自成进入河南}

\eventanchor{ML-1641-002}{1641（崇祯十四年）三月·李自成攻克洛阳}

\eventanchor{ML-1641-003}{1641（崇祯十四年）·起义军首领张献忠攻克襄阳}

\eventanchor{ML-1641-004}{1641（崇祯十四年）七月15日·张献忠拿下武昌}

\eventanchor{ML-1641-005}{1641（崇祯十四年）·蝗虫袭击浙江}

\eventanchor{ML-1641-006}{1641（崇祯十四年）·辽东防务、后金（清）军事行动和流动作战集团的本年记录标记多战场互动，军数和胜败须保留分歧。}

\eventanchor{ML-1641-007}{1641（崇祯十四年）·地方结社、宗教、出版或士人文集的本年线索提供战乱中的社会观察，幸存者记忆不能概括全国。}

\subsection{事件导航与来源边界}

\eventanchor{ML-1642-001}{1642（崇祯十五年）四月8日·宋金之战：清军占领锦州}

\eventanchor{ML-1642-002}{1642（崇祯十五年）·浙江遭遇洪涝灾害}

\eventanchor{ML-1642-003}{1642（崇祯十五年）·复合金属大炮产生于明代。}

\eventanchor{ML-1642-004}{1642（崇祯十五年）·李自成的叛军成功地用大炮在明朝的防御工事上造成了两丈的缺口。}

\eventanchor{ML-1642-005}{1642（崇祯十五年）·辽饷、剿饷、练饷及地方征解的本年文书显示财政负担，定额、征收与实际流向不得混为一谈。}

\eventanchor{ML-1642-006}{1642（崇祯十五年）·朝鲜、辽东和海上关系的本年材料提供外部观察位置，明、后金（清）与地方势力的称谓需具体化。}

\eventanchor{ML-1642-007}{1642（崇祯十五年）·旱灾、蝗灾、疫病、河患与赈济的本年材料连接环境和社会压力，死亡与流离数字须谨慎处理。}

\eventanchor{ML-1643-001}{1643（崇祯十六年）一月·李自成攻克襄阳}

\eventanchor{ML-1643-002}{1643（崇祯十六年）十一月·李自成攻克西安}

\eventanchor{ML-1643-003}{1643（崇祯十六年）·张献忠在湖广宣告奚朝}

\eventanchor{ML-1643-004}{1643（崇祯十六年）·辽东防务、后金（清）军事行动和流动作战集团的本年记录标记多战场互动，军数和胜败须保留分歧。（材料核查重点：诏令或奏议的形成与发受链）}

\eventanchor{ML-1643-005}{1643（崇祯十六年）·辽东防务、后金（清）军事行动和流动作战集团的本年记录标记多战场互动，军数和胜败须保留分歧。（材料核查重点：报称信息的来源、抄传与行政分类）}

\eventanchor{ML-1643-006}{1643（崇祯十六年）·地方结社、宗教、出版或士人文集的本年线索提供战乱中的社会观察，幸存者记忆不能概括全国。}

\eventanchor{ML-1643-007}{1643（崇祯十六年）·天启末至崇祯朝的任免、奏疏和廷议构成本年中枢危机线索；崇祯无实录，必须以多类材料互校。}

\eventanchor{ML-1644-001}{1644（崇祯十七年）二月8日·李自成在西安宣告大顺}

\eventanchor{ML-1644-002}{1644（崇祯十七年）四月25日·李自成占领北京，崇祯皇帝上吊自杀}

\eventanchor{ML-1644-003}{1644（崇祯十七年）三月至四月·京畿防务、流动作战与军饷决策的文书显示崇祯十七年初的中枢危机；各路兵力和命令传递需要多源校验。}

\eventanchor{ML-1644-004}{1644（崇祯十七年）·朝鲜、辽东和海上关系的本年材料提供外部观察位置，明、后金（清）与地方势力的称谓需具体化。（材料核查重点：诏令或奏议的形成与发受链）}

\eventanchor{ML-1644-005}{1644（崇祯十七年）·朝鲜、辽东和海上关系的本年材料提供外部观察位置，明、后金（清）与地方势力的称谓需具体化。（材料核查重点：报称信息的来源、抄传与行政分类）}

\eventanchor{ML-1644-006}{1644（崇祯十七年）·旱灾、蝗灾、疫病、河患与赈济的本年材料连接环境和社会压力，死亡与流离数字须谨慎处理。}

\eventanchor{ML-1644-007}{1644（崇祯十七年）·辽饷、剿饷、练饷及地方征解的本年文书显示财政负担，定额、征收与实际流向不得混为一谈。}
